\documentclass[../../main.tex]{subfiles}

\begin{document}

\section{度量空间中的紧致性}\label{section:ch07-compactness-metric}

我们已经证明了对于度量空间来说, 紧致性、极限点紧致性以及列紧性都是等价的. 度量空间中的紧致性也有其他的描述方法, 其中有一种就涉及到完备性的概念. 这一节我们就来研究它. 作为应用, 我们再证明一个刻划 $\mathcal{C}(X,\mathbb{R}^n)$ 中 (关于一致度量) 紧致子集的定理.

度量空间 $X$ 的紧致性与它的完备性有什么关系呢? 根据\cref{lemma:ch07-cauchy-subsequence}可见, 每一个紧致度量空间必定是完备的. 反之未必成立, 也就是说, 完备度量空间未必是紧致的. 那么下面一个很自然的问题就是, 完备空间需附加什么条件才能得出紧致性呢? 这个条件就是所谓的完全有界性.

\begin{definition}\label{definition:ch07-totally-bounded}
度量空间 $(X,d)$ 称为完全有界的(totally bounded), 如果对任意 $\varepsilon>0$, 存在一个由 $\varepsilon$-球构成的 $X$ 的有限覆盖.
\end{definition}

\begin{example}\label{example:ch07-totally-bounded-vs-bounded}
显然, 度量空间的完全有界性蕴涵有界性. 因为若 $B(x_1,1/2),\cdots,B(x_n,1/2)$ 是 $X$ 的一个由半径为 $1/2$ 的开球构成的有限覆盖, 则 $X$ 的直径最多为 $1+\max\{d(x_i,x_j)\}$. 但反之未必成立. 例如, 在度量 $\bar{d}(a,b)=\min\{1,|a-b|\}$ 之下, 实直线 $\mathbb{R}$ 是有界的但不是完全有界的.
\end{example}

\begin{example}\label{example:ch07-complete-totally-bounded}
关于度量 $d(a,b)=|a-b|$, 实直线 $\mathbb{R}$ 是完备的但不是完全有界的. 子空间 $(-1,1)$ 是完全有界的, 却不是完备的, 而子空间 $[-1,1]$ 既是完备的, 也是完全有界的.
\end{example}

\begin{theorem}\label{theorem:ch07-compact-complete-totally}
度量空间 $(X,d)$ 是紧致的, 当且仅当它是完备和完全有界的.
\end{theorem}

\begin{proof}
如上所说, 若 $X$ 是一个紧致度量空间, 则 $X$ 自然是完备的. 又因为所有的开 $\varepsilon$-球做成的 $X$ 的开覆盖必包含一个有限子覆盖, 所以 $X$ 是完全有界的.

反之, 设 $X$ 是完备的和完全有界的. 下面只需证明 $X$ 是列紧的便可以了.

设 $(x_n)$ 是 $X$ 中点的一个序列, 我们来选取它的一个子序列, 要求这个子序列是一个 Cauchy 序列, 因而一定收敛. 首先, 用有限多个半径为 $1$ 的球覆盖 $X$, 这时, 至少有一个球, 比如说 $B_1$, 包含无限多个 $x_n$. 设 $J_1$ 是 $\mathbb{Z}_+$ 的一个子集, 使得对于每一个 $x_n$, 若 $x_n\in B_1$, 则 $n\in J_1$.

其次, 用有限多个半径为 $1/2$ 的球覆盖 $X$, 因为 $J_1$ 是无限的, 这些球中至少有一个, 比如说 $B_2$, 包含 $J_1$ 中无限多个 $n$. 取 $J_2$ 为满足 $n\in J_1$ 和 $x_n\in B_2$ 的所有 $n$ 构成的集合. 一般地, 给定正整数的无穷集 $J_k$, 取 $J_{k+1}$ 为 $J_k$ 的一个无穷子集, 使得存在一个半径为 $1/(k+1)$ 的球 $B_{k+1}$, 对于所有的 $n\in J_{k+1}$, 有 $x_n\in B_{k+1}$.

选取 $n_1\in J_1$. 对于给定的 $n_k$, 选取 $n_{k+1}\in J_{k+1}$, 使得 $n_{k+1}>n_k$. 因为 $J_{k+1}$ 是无穷集, 这一点是可以做得到的. 现在, 对于 $i,j\geqslant k$, 指标 $n_i,n_j$ 都属于 $J_k$ (因为 $J_1\supset J_2\supset\cdots$ 是一个集合的套序列). 因此, 对于所有的 $i,j\geqslant k$, 点 $x_{n_i}$ 和 $x_{n_j}$ 都包含在半径为 $1/k$ 的球 $B_k$ 中. 于是, 序列 $(x_{n_i})$ 是一个 Cauchy 序列.
\end{proof}

我们现在应用这一结论来寻找 $\mathcal{C}(X,\mathbb{R}^n)$ 在一致拓扑下的一个紧致子空间. 我们已经知道 $\mathbb{R}^n$ 的子空间是紧致的当且仅当它是闭的和有界的. 大家可能希望对于 $\mathcal{C}(X,\mathbb{R}^n)$ 也有相似的结论成立. 但事实并非如此, 即使在 $X$ 紧致的情况下也不行. 其实, 这需要 $\mathcal{C}(X,\mathbb{R}^n)$ 的子空间满足一个条件, 即等度连续性. 现在我们来看它的定义.

\begin{definition}\label{definition:ch07-equicontinuous}
设 $(Y,d)$ 是一个度量空间. $\mathcal{F}$ 是函数空间 $\mathcal{C}(X,Y)$ 的一个子集. 对于 $x_0\in X$, 函数集 $\mathcal{F}$ 称为在 $x_0$ 处等度连续的(equicontinuous at $x_0$), 如果给定 $\varepsilon>0$, 存在 $x_0$ 的一个邻域 $U$, 使得对于所有的 $x\in U$ 及所有 $f\in\mathcal{F}$ 有
\begin{align*}
d(f(x),f(x_0))<\varepsilon.
\end{align*}
如果对于每一个点 $x_0\in X$, 集合 $\mathcal{F}$ 在 $x_0$ 处等度连续, 则称 $\mathcal{F}$ 是等度连续的(equicontinuous).
\end{definition}

函数 $f$ 在点 $x_0$ 处的连续性是指对给定的 $f$ 及给定的 $\varepsilon>0$, 存在 $x_0$ 的一个邻域 $U$, 使对于所有 $x\in U$ 有 $d(f(x),f(x_0))<\varepsilon$. $\mathcal{F}$ 的等度连续性是指存在一个邻域 $U$, 对于 $\mathcal{F}$ 中的所有的函数 $f$, 上式都成立.

注意, 等度连续性依赖于具体的度量, 而不仅仅是 $Y$ 上的拓扑.

\begin{lemma}\label{lemma:ch07-totally-bounded-equicontinuous}
设 $X$ 是一个空间, $(Y,d)$ 是一个度量空间. 设 $\mathcal{F}$ 是 $\mathcal{C}(X,Y)$ 的一个子集, 并且关于度量 $d$ 的一致度量是完全有界的, 则 $\mathcal{F}$ 关于 $d$ 是等度连续的.
\end{lemma}

\begin{proof}
设 $\mathcal{F}$ 是完全有界的. 给定 $0<\varepsilon<1$, 给定 $x_0$, 以下证明存在 $x_0$ 的一个邻域 $U$, 对于所有 $f\in\mathcal{F}$ 和 $x\in U$, 有 $d(f(x),f(x_0))<\varepsilon$.

设 $\delta=\varepsilon/3$, 用 $\mathcal{C}(X,Y)$ 中的有限多个开 $\delta$-球
\begin{align*}
B(f_1,\delta),\cdots,B(f_n,\delta)
\end{align*}
覆盖 $\mathcal{F}$. 因为每一个函数 $f_i$ 是连续的, 我们可以选取 $x_0$ 的一个邻域 $U$, 使得对于 $i=1,\cdots,n$ 和 $x\in U$, 有
\begin{align*}
d(f_i(x),f_i(x_0))<\delta.
\end{align*}

设 $f$ 是 $\mathcal{F}$ 的任意一个元素, 则 $f$ 至少属于上述 $\delta$-球中的某一个, 不妨设为 $B(f_i,\delta)$. 于是对于 $x\in U$, 有
\begin{align*}
\bar{d}(f(x),f_i(x))<\delta,\\
d(f_i(x),f_i(x_0))<\delta,\\
\bar{d}(f_i(x_0),f(x_0))<\delta.
\end{align*}
第一个和第三个不等式成立是因为 $\bar{\rho}(f,f_i)<\delta$, 第二个不等式成立是因为 $x\in U$. 由于 $\delta<1$, 所以用 $d$ 代替 $\bar{d}$ 时, 第一个和第三个不等式依然成立. 所以根据三角不等式可见对于所有的 $x\in U$, $d(f(x),f(x_0))<\varepsilon$ 成立.
\end{proof}

现在我们来证明 Ascoli 定理的经典形式, 它所讨论的就是函数空间 $\mathcal{C}(X,\mathbb{R}^n)$ 的紧致子空间. 这个定理的一个更广义形式, 其证明并不依赖这个经典形式 (将在第 47 节中给出). 广义形式的证明依赖于 Tychonoff 定理, 但在这里并不需要.

\begin{lemma}\label{lemma:ch07-equicontinuous-totally-bounded}
设 $X$ 是一个空间, $(Y,d)$ 是一个度量空间. 设 $X$ 和 $Y$ 是紧致的. 如果 $\mathcal{F}$ 关于 $d$ 是 $\mathcal{C}(X,Y)$ 的一个等度连续子集, 那么关于度量 $d$ 的一致度量和上确界度量, $\mathcal{F}$ 是完全有界的.
\end{lemma}

\begin{proof}
由于 $X$ 是紧致的, 可以在 $\mathcal{C}(X,Y)$ 上定义上确界度量 $\rho$. 关于 $\rho$ 完全有界等价于关于 $\bar{\rho}$ 完全有界. 因为对于 $\varepsilon<1$, 关于 $\rho$ 的每一个 $\varepsilon$-球也是关于 $\bar{\rho}$ 的一个 $\varepsilon$-球, 反之亦成立. 所以我们可以始终只用度量 $\rho$.

设 $\mathcal{F}$ 是等度连续的. 给定 $\varepsilon>0$, 用关于度量 $\rho$ 的有限多个开 $\varepsilon$-球覆盖 $\mathcal{F}$.

令 $\delta=\varepsilon/3$. 任意给定 $a\in X$, 存在 $a$ 的相应邻域 $U_a$ 使得对任意 $x\in U_a$ 和任意 $f\in\mathcal{F}$, 有
\begin{align*}
d(f(x),f(a))<\delta.
\end{align*}
于是可以找到有限多个邻域 $U_a$, $a=a_1,\cdots,a_k$, 组成 $X$ 的一个覆盖, 用 $U_i$ 表示 $U_{a_i}$. 选取直径小于 $\delta$ 的有限个开集 $V_1,\cdots,V_m$ 覆盖 $Y$.

设 $J$ 是所有函数 $\alpha\colon\{1,\cdots,k\}\to\{1,\cdots,m\}$ 的集合. 给定 $\alpha\in J$, 若存在 $\mathcal{F}$ 中的一个函数 $f$, 使得对于每一个 $i=1,\cdots,k$, $f(a_i)\in V_{\alpha(i)}$, 便选取一个这样的函数, 记为 $f_\alpha$. 集合 $\{f_\alpha\}$ 是由集合 $J$ 的一个子集 $J'$ 加标的, 所以是有限的. 我们断言: 开球 $B_\rho(f_\alpha,\varepsilon)$, $\alpha\in J'$, 覆盖 $\mathcal{F}$.

设 $f$ 是 $\mathcal{F}$ 的一个元素. 对于每一个 $i=1,\cdots,k$, 选取整数 $\alpha(i)$ 使得 $f(a_i)\in V_{\alpha(i)}$. 则函数 $\alpha$ 在 $J'$ 中. 我们断言: $f$ 属于球 $B_\rho(f_\alpha,\varepsilon)$.

设 $x$ 是 $X$ 的一个点. 选取 $i$ 使得 $x\in U_i$. 于是
\begin{align*}
d(f(x),f(a_i))<\delta,\\
d(f(a_i),f_\alpha(a_i))<\delta,\\
d(f_\alpha(a_i),f_\alpha(x))<\delta.
\end{align*}
第一个和第三个不等式成立是因为 $x\in U_i$, 第二个不等式成立是因为 $f(a_i)$ 和 $f_\alpha(a_i)$ 都属于 $V_{\alpha(i)}$. 综上可见, $d(f(x),f_\alpha(x))<\varepsilon$. 因对于所有的 $x\in X$, 以上不等式成立, 所以
\begin{align*}
\rho(f,f_\alpha)=\max\{d(f(x),f_\alpha(x))\}<\varepsilon.
\end{align*}
因此 $f$ 属于开球 $B_\rho(f_\alpha,\varepsilon)$.
\end{proof}

\begin{definition}\label{definition:ch07-pointwise-bounded}
设 $(Y,d)$ 是一个度量空间. $\mathcal{C}(X,Y)$ 的子集 $\mathcal{F}$ 称为关于 $d$ 是点态有界的(pointwise bounded), 如果对于每一个 $a\in X$, $Y$ 的子集
\begin{align*}
\mathcal{F}_a=\{f(a)\mid f\in\mathcal{F}\}
\end{align*}
关于 $d$ 是有界的.
\end{definition}

\begin{theorem}[Ascoli 定理, 经典形式]\label{theorem:ch07-ascoli-classical}
设 $X$ 是一个紧致空间, $(\mathbb{R}^n,d)$ 表示关于平方度量或欧氏度量的欧氏空间, 赋予 $\mathcal{C}(X,\mathbb{R}^n)$ 相应的一致拓扑. 那么 $\mathcal{C}(X,\mathbb{R}^n)$ 的子空间 $\mathcal{F}$ 有紧致闭包当且仅当 $\mathcal{F}$ 关于 $d$ 是等度连续的和点态有界的.
\end{theorem}

\begin{proof}
因为 $X$ 是紧致的, 在 $\mathcal{C}(X,\mathbb{R}^n)$ 上定义上确界度量 $\rho$, 并赋予 $\mathcal{C}(X,\mathbb{R}^n)$ 一致拓扑. 并始终用 $\mathcal{G}$ 表示 $\mathcal{F}$ 在 $\mathcal{C}(X,\mathbb{R}^n)$ 中的闭包.

第一步. 证明若 $\mathcal{G}$ 是紧致的, 则关于 $d$, $\mathcal{G}$ 是等度连续的并且是点态有界的. 根据 $\mathcal{F}\subset\mathcal{G}$, 可见 $\mathcal{F}$ 也是等度连续的并且点态有界的. 这就证明了定理的"必要性".

根据\cref{theorem:ch07-compact-complete-totally}可见, $\mathcal{G}$ 的紧致性蕴涵着 $\mathcal{G}$ 关于 $\rho$ 和关于 $\bar{\rho}$ 都是完全有界的, 再根据\cref{lemma:ch07-totally-bounded-equicontinuous}, $\mathcal{G}$ 关于 $d$ 是等度连续的. $\mathcal{G}$ 的紧致还蕴涵着 $\mathcal{G}$ 关于 $\rho$ 是有界的. 由此又可见 $\mathcal{G}$ 在 $d$ 下是点态有界的. 因为如果对于任意 $f,g\in\mathcal{G}$ 有 $\rho(f,g)\leqslant M$, 那么, 特别地, $d(f(a),g(a))\leqslant M$, 因此 $\mathcal{G}_a$ 的直径最多为 $M$.

第二步. 证明若 $\mathcal{F}$ 关于 $d$ 是等度连续和点态有界的, 则 $\mathcal{G}$ 也是等度连续和点态有界的.

首先验证等度连续性. 给定 $x_0\in X$ 和 $\varepsilon>0$, 选取 $x_0$ 的一个邻域 $U$, 使得对于所有的 $x\in U$ 和 $f\in\mathcal{F}$, 有 $d(f(x),f(x_0))<\varepsilon/3$. 任意给定 $g\in\mathcal{G}$, 选取 $f\in\mathcal{F}$, 使得 $\rho(f,g)<\varepsilon/3$. 根据三角不等式可见, 对于所有 $x\in U$, 有 $d(g(x),g(x_0))<\varepsilon$. 因为 $g$ 是任意的, 所以 $\mathcal{G}$ 在 $x_0$ 点等度连续.

其次再来证明点态有界性. 给定 $a$, 选取 $M$, 使得 $\mathcal{F}_a$ 的直径不大于 $M$. 任意给定 $g,g'\in\mathcal{G}$, 选取 $f,f'\in\mathcal{F}$, 使得 $\rho(f,g)<1$ 和 $\rho(f',g')<1$. 因为 $d(f(a),f'(a))\leqslant M$, 所以 $d(g(a),g'(a))\leqslant M+2$. 再由 $g$ 和 $g'$ 是任意的可得 $\operatorname{diam}\mathcal{G}_a\leqslant M+2$.

第三步. 证明若 $\mathcal{G}$ 是等度连续的并且点态有界的, 则存在 $\mathbb{R}^n$ 的紧致子空间 $Y$, 它包含着所有集合 $g(X)$ 的并, 其中 $g\in\mathcal{G}$.

对于每一个 $a\in X$, 选取邻域 $U_a$, 使得对于所有 $x\in U_a$ 和 $g\in\mathcal{G}$, 有 $d(g(x),g(a))<1$. 因为 $X$ 是紧致的, 所以可以选取有限多个这种邻域覆盖 $X$, 不妨设这些邻域的下标为 $a=a_1,\cdots,a_k$. 因为 $\mathcal{G}_{a_i}$ 是有界的, 所以它们的并也是有界的, 设它包含于 $\mathbb{R}^n$ 中的以原点为球心、以 $N$ 为半径的球中. 于是对于所有的 $g\in\mathcal{G}$, 集合 $g(X)$ 包含于以原点为球心、以 $N+1$ 为半径的球中. 令 $Y$ 为这个球的闭包.

第四步. 充分性的证明. 设 $\mathcal{F}$ 关于 $d$ 是等度连续和点态有界的. 我们证明关于 $\rho$, $\mathcal{G}$ 是完备的和完全有界的. 根据\cref{theorem:ch07-compact-complete-totally}可见, $\mathcal{G}$ 是紧致的.

完备性容易证得. 因为 $\mathcal{G}$ 是完备度量空间 $(\mathcal{C}(X,\mathbb{R}^n),\rho)$ 的一个闭子空间.

我们来验证完全有界性. 首先, 根据第二步可见 $\mathcal{G}$ 在度量 $d$ 下是等度连续的和点态有界的; 再根据第三步可见, 存在 $\mathbb{R}^n$ 的紧致子空间 $Y$, 使得 $\mathcal{G}\subset\mathcal{C}(X,Y)$. 根据\cref{lemma:ch07-equicontinuous-totally-bounded}, $\mathcal{G}$ 的等度连续性蕴涵着 $\mathcal{G}$ 关于 $\rho$ 是完全有界的.
\end{proof}

\begin{corollary}\label{corollary:ch07-ascoli-classical}
设 $X$ 是紧致的, $d$ 为 $\mathbb{R}^n$ 上的平方度量或欧氏度量, 赋予 $\mathcal{C}(X,\mathbb{R}^n)$ 相应的一致拓扑. 则 $\mathcal{C}(X,\mathbb{R}^n)$ 的一个子空间 $\mathcal{F}$ 是紧致的, 当且仅当它关于上确界度量 $\rho$ 是闭的和有界的, 关于 $d$ 是等度连续的.
\end{corollary}

\begin{proof}
若 $\mathcal{F}$ 是紧致的, 则它必定是闭的和有界的. 根据前一个定理易见它还是等度连续的. 反之, 若 $\mathcal{F}$ 是闭的, 则它就等于它的闭包 $\mathcal{G}$. 若它关于 $\rho$ 是有界的, 则关于 $d$ 是点态有界的. 若它是等度连续的, 则前一个定理蕴涵着它是紧致的.
\end{proof}

\subsection*{习题}

\begin{enumerate}
\item 若 $X_n$ 是可度量化的, 并且有度量 $d_n$, 则
\begin{align*}
D(\mathbf{x},\mathbf{y})=\sup\{\bar{d}_i(x_i,y_i)/i\}
\end{align*}
是积空间 $X=\prod X_n$ 的一个度量. 证明若每一个 $X_n$ 在 $d_n$ 下都是完全有界的, 则 $X$ 在 $D$ 下是完全有界的. 在不使用 Tychonoff 定理的情况下, 证明紧致可度量化空间的可数积空间是紧致的.
\item 设 $(Y,d)$ 是一个度量空间, $\mathcal{F}$ 是 $\mathcal{C}(X,Y)$ 的一个子集.
\begin{enumerate}[(a)]
\item 证明: 若 $\mathcal{F}$ 是有限的, 则 $\mathcal{F}$ 是等度连续的.
\item 证明: 若 $f_n$ 是 $\mathcal{C}(X,Y)$ 中元素的一个序列, 并且一致收敛, 则集合 $\{f_n\}$ 是等度连续的.
\item 设 $\mathcal{F}$ 是可微函数 $f\colon\mathbb{R}\to\mathbb{R}$ 的集合, 其中这些可微函数满足条件: 对于每一个 $x\in\mathbb{R}$, 存在某一个邻域 $U$, 使得 $\mathcal{F}$ 的所有函数的导数在这个邻域 $U$ 上是一致有界的. [即存在 $M$, 使得对于所有 $y\in U$ 和所有 $f\in\mathcal{F}$ 有 $|f'(y)|\leqslant M$.] 证明 $\mathcal{F}$ 是等度连续的.
\end{enumerate}
\item 证明:
\begin{quote}
\textbf{定理(Arzelà 定理).} 设 $X$ 紧致的, 并且 $f_n\in\mathcal{C}(X,\mathbb{R}^k)$. 若族 $\{f_n\}$ 是点态有界的和等度连续的, 则序列 $f_n$ 有一个一致收敛的子序列.
\end{quote}
\item (a) 设函数 $f_n\colon I\to\mathbb{R}$ 定义为 $f_n(x)=x^n$. 集合 $\mathcal{F}=\{f_n\}$ 是点态有界的, 但序列 $(f_n)$ 没有一致收敛的子序列, 在哪个或哪些点处 $\mathcal{F}$ 不是等度连续的?
\item (b) 针对第 21 节的习题 9 中定义的函数 $f_n$ 再次讨论 (a) 的情况.
\item 设 $X$ 是一个空间. $\mathcal{C}(X,\mathbb{R})$ 的子集 $\mathcal{F}$ 称为在无穷远处一致蜕化(vanish uniformly at infinity), 如果给定 $\varepsilon>0$, 存在 $X$ 的一个紧致子集 $C$, 使得对于每一个 $x\in X-C$ 及 $f\in\mathcal{F}$, 有 $|f(x)|<\varepsilon$. 如果 $\mathcal{F}$ 仅由一个函数 $f$ 组成, 则称 $f$ 在无穷远处蜕化(vanishes at infinity). 用 $\mathcal{C}_0(X,\mathbb{R})$ 表示在无穷远处蜕化的连续函数 $f\colon X\to\mathbb{R}$ 的族.
\begin{quote}
\textbf{定理.} 设 $X$ 是局部紧致的 Hausdorff 空间, 赋予 $\mathcal{C}_0(X,\mathbb{R})$ 一致拓扑. $\mathcal{C}_0(X,\mathbb{R})$ 的子集 $\mathcal{F}$ 有紧致闭包, 当且仅当它是点态有界的和等度连续的, 并且在无穷远处一致蜕化.
\end{quote}
[提示: 用 $Y$ 表示 $X$ 的单点紧致化, 证明若赋予 $\mathcal{C}_0(X,\mathbb{R})$ 和 $\mathcal{C}(Y,\mathbb{R})$ 上确界度量, 则 $\mathcal{C}_0(X,\mathbb{R})$ 等距于 $\mathcal{C}(Y,\mathbb{R})$ 的一个闭子空间.]
\item 在 Ascoli 定理的证明中将 $\mathbb{R}^n$ 统统换为任意一个度量空间, 只要求它满足条件: 其中的所有闭的有界子空间都是紧致的. 证明结论依然成立.
\item 设 $(X,d)$ 是一个度量空间. 如果 $A\subset X$ 和 $\varepsilon>0$, 令 $U(A,\varepsilon)$ 表示 $A$ 的 $\varepsilon$-邻域. 令 $\mathcal{H}$ 为 $X$ 的所有 (非空) 有界闭子集的族. 如果 $A,B\in\mathcal{H}$, 定义
\begin{align*}
D(A,B)=\inf\{\varepsilon\mid A\subset U(B,\varepsilon)\text{ 并且 }B\subset U(A,\varepsilon)\}.
\end{align*}
\begin{enumerate}[(a)]
\item 证明: $D$ 是 $\mathcal{H}$ 的一个度量, 称之为 Hausdorff 度量(Hausdorff metric).
\item 证明: 如果 $(X,d)$ 是完备的, 则 $(\mathcal{H},D)$ 也是完备的. [提示: 设 $A_n$ 为 $\mathcal{H}$ 中的一个 Cauchy 序列. 用过渡到子序列的办法, 假定 $D(A_n,A_{n+1})<1/2^n$. 定义 $A$ 为序列 $x_1,x_2,\cdots$ 的极限点的集合, 要求这些序列 $x_1,x_2,\cdots$ 满足条件: 对于每一个 $i$, $x_i\in A_i$, 并且 $d(x_i,x_{i+1})<1/2^i$. 证明 $A_n\to\overline{A}$.]
\item 证明: 若 $(X,d)$ 是完全有界的, 则 $(\mathcal{H},D)$ 也是完全有界的. [提示: 给定 $\varepsilon$, 选取 $\delta<\varepsilon$, 并且令 $S$ 是 $X$ 的一个有限子集, 使得族 $\{B_d(x,\delta)\mid x\in S\}$ 覆盖 $X$. 设 $\mathcal{A}$ 是 $S$ 的所有非空子集构成的族. 证明 $\{B_D(A,\varepsilon)\mid A\in\mathcal{A}\}$ 覆盖 $\mathcal{H}$.]
\item \textbf{定理.} 若 $X$ 关于度量 $d$ 是紧致的, 则 $\mathcal{H}$ 关于 Hausdorff 度量 $D$ 也是紧致的.
\end{enumerate}
\item 设 $(X,d_X)$ 和 $(Y,d_Y)$ 都是度量空间, 赋予 $X\times Y$ 相应的平方度量, 用 $\mathcal{H}$ 表示 $X\times Y$ 关于 Hausdorff 度量 $D$ 的所有 (非空) 有界闭子集的族. 关于一致度量, 考虑空间 $\mathcal{C}(X,Y)$. 设映射 $gr\colon\mathcal{C}(X,Y)\to\mathcal{H}$ 是在每一个连续函数 $f\colon X\to Y$ 上, 取值为函数 $f$ 的图像
\begin{align*}
G_f=\{x\times f(x)\mid x\in X\}
\end{align*}
的那个函数.
\begin{enumerate}[(a)]
\item 证明映射 $gr$ 是单射, 并且一致连续.
\item 用 $\mathcal{H}_0$ 表示映射 $gr$ 的像集, 设 $g\colon\mathcal{C}(X,Y)\to\mathcal{H}_0$ 是由 $gr$ 得到的一个满射. 证明若 $f\colon X\to Y$ 是一致连续的, 则映射 $g^{-1}$ 在点 $G_f$ 处是连续的.
\item 举一个映射 $g^{-1}$ 在点 $G_f$ 处不连续的例子.
\item \textbf{定理.} 若 $X$ 是紧致的, 则 $gr\colon\mathcal{C}(X,Y)\to\mathcal{H}$ 是一个嵌入.
\end{enumerate}
\end{enumerate}

\end{document}